\title{Direct Preference Optimization: Your Language Model is Secretly a Reward Model}

\begin{abstract}
Large-scale unsupervised language models (LMs) acquire extensive world knowledge and certain reasoning abilities, yet directing their behavior with precision remains challenging because their training is entirely unsupervised. Traditional approaches to enhance model steerability involve collecting human judgements on the comparative quality of model outputs and adjusting the unsupervised LM accordingly—typically through reinforcement learning from human feedback (RLHF). However, RLHF often entails a complicated and sometimes unstable two-step process: first constructing a reward model that captures human preferences, and then applying reinforcement learning to fine-tune the base LM so as to optimize this proxy reward, all while constraining the model to remain close to its initial state. In this work, we propose a novel reparameterization of the reward model within RLHF that allows for direct, closed-form computation of the optimal policy. This reformulation reduces the RLHF optimization to a straightforward classification loss. We introduce the resulting technique as \textit{Direct Preference Optimization} (DPO), which demonstrates strong, reliable performance while also being computationally efficient—removing the need for LM sampling during training or extensive hyperparameter searches. Experimental results indicate that DPO fine-tunes LMs to human preferences at least as effectively as, and often better than, established techniques. In particular, DPO outperforms PPO-based RLHF in sentiment control, and matches or surpasses its performance in tasks such as summarization and single-turn dialogue, while maintaining a significantly simpler training procedure and implementation.
\end{abstract}

\section{Introduction}
Large-scale unsupervised language models (LMs) that are trained on massive corpora have demonstrated a wide range of unexpected abilities~\citep{chowdhery2022palm, brown2020language, touvron2023llama,bubeck2023sparks}. Nonetheless, these models are exposed to data generated by humans with diverse objectives, levels of expertise, and intentions. Not all of these human-generated behaviors are appropriate to replicate; for instance, it is desirable for an AI code assistant to \textit{recognize} prevalent programming errors to rectify them, yet when producing code, we prefer the model to reflect the higher-caliber—albeit less common—coding practices available in its training corpus. In a similar vein, while a language model should be \textit{informed} of a widely held misconception (e.g., one believed by 50\% of people), we clearly do not want the model to reproduce this misconception as fact in half of the relevant queries. Thus, it becomes essential to differentiate the model’s \emph{intended outputs and behaviors} from the broader array of \textit{knowledge and capabilities} it acquires, so that AI systems remain reliable, effective, and amenable to control \citep{ouyang2022training}. Current techniques typically adjust LMs to align with human preferences through reinforcement learning (RL), but we argue that the objective function used in these RL-based approaches can, in fact, be directly minimized via a simple binary cross-entropy objective, thereby streamlining the overall preference learning workflow.

Conceptually, the predominant strategies impart desirable conduct to a language model by leveraging human-curated preference data that embodies traits considered safe and constructive. This phase of learning from preferences takes place after an initial, broad unsupervised pre-training stage on extensive textual resources. A straightforward path to aligning with preferences is to apply supervised fine-tuning on demonstrations of preferred responses. Nevertheless, the most influential and effective category of approaches—namely, reinforcement learning from human or AI feedback (RLHF/RLAIF; \citep{christiano2017deep,bai2022constitutional})—instead entails training a reward model on annotated human preferences, which is then utilized to optimize the LM policy via RL so that responses with higher reward become more probable, all while preventing substantial divergence from the pretrained model. While RLHF-based techniques have yielded models with exceptional abilities in dialogue and coding, the RLHF pipeline is markedly more intricate than traditional supervised training, demanding the development of several models and iterative sampling from the language model policy during training, leading to notable computational overhead.

This paper presents a method for directly optimizing language models to align with human preferences, omitting the need for explicit reward modeling or reinforcement learning. We introduce \textit{Direct Preference Optimization} (DPO), an algorithm that effectively targets the same objective as RLHF methods—namely, maximizing reward under a KL-divergence constraint—but is significantly easier to implement and train. Conceptually, DPO adjusts the model by boosting the relative log-likelihood of responses deemed preferable, as compared to those regarded as less preferable. It achieves this by integrating a dynamic importance weighting for each data instance, thereby mitigating the model collapse observed with a simplistic probability ratio formulation. Like prior approaches, DPO is grounded in a theoretical framework for preferences (for instance, the Bradley-Terry model; \cite{bradley1952rankanalysis}), which assesses the correspondence between a reward function and observed preference judgments. In contrast to conventional strategies—where the preference model informs the loss for training a separate reward function, followed by optimizing a policy with respect to that learned reward—DPO reparameterizes the preference loss, expressing it directly in terms of the policy parameters. As a result, with access to human-labeled preference data for model outputs, DPO can optimize the model using a straightforward binary cross-entropy objective, producing a policy that implicitly maximizes an unknown but best-fitting reward function derived from the empirical preferences.

Our principal contribution is Direct Preference Optimization (DPO): a streamlined, reinforcement learning-free technique for preference-based language model training. Empirical results demonstrate that DPO matches or exceeds the effectiveness of current techniques—including RLHF implementations based on PPO—across preference-driven tasks like sentiment control, text summarization, and dialogue modeling, in language models scaling up to 6B parameters.

\section{Related Work}

Large self-supervised language models demonstrate the capability to handle certain tasks in a zero-shot manner \citep{radford2019language} or through few-shot prompting \citep{gpt3,megatron,chowdhery2022palm}. Nevertheless, their effectiveness on downstream applications and alignment with user objectives are markedly enhanced when they undergo fine-tuning using instruction-based datasets paired with human-authored completions \citep{mishra-etal-2022-cross,sanh2022multitask,chung2022scaling,thoppilan2022lamda}. Through the process known as `instruction-tuning,' LLMs gain the ability to generalize beyond their initial training instructions, ultimately boosting their practical utility \citep{chung2022scaling}. Although instruction tuning yields substantial improvements, it is often more feasible to amass \textit{relative} human assessments of output quality than expert-level reference completions. Consequently, subsequent research has adopted the practice of further fine-tuning LLMs using collections of human preferences, yielding gains in areas such as translation \citep{kreutzer-etal-2018-reliability}, summarization \citep{stiennon2022learning,ziegler2020finetuning}, narrative generation \citep{ziegler2020finetuning}, and adherence to instructions \citep{ouyang2022training,ramamurthy2023is}. This approach typically entails first training a neural reward model aligned with the preference data according to a framework such as the Bradley-Terry model \citep{bradley1952rankanalysis}, and then adjusting the language model's outputs to optimize for these rewards through reinforcement learning strategies like REINFORCE \citep{williams1992reinforce}, proximal policy optimization (PPO; \cite{schulman2017proximal}), or their adaptations \citep{ramamurthy2023is}. Another related strategy employs LLMs already instruction-tuned via human feedback to produce further synthetic preference data, particularly for desired qualities like safety or harmlessness \citep{bai2022constitutional}, relying mainly on weak supervision using textual rubrics provided by humans for guiding the model's responses. These developments illustrate the merging of two streams of research: reinforcement learning-based training of language models for assorted objectives~\citep{Ranzato2015SequenceLT,paulus2018a,wu2018learning}, and methods geared toward extracting knowledge from human preferences \citep{christiano2017deep,kupcsik2018learning}. Yet, while employing relative preferences is intuitively advantageous, employing reinforcement learning for fine-tuning large language models poses substantial practical obstacles. The present work introduces a principled method for directly optimizing with respect to relative preferences, circumventing the need for reinforcement learning.

Beyond linguistic applications, the problem of learning policies from preference data has been explored in both bandit and reinforcement learning contexts, resulting in a range of proposed methods. In particular, the contextual dueling bandit (CDB; \cite{yue2012karmed,dudik2015contextual}) framework considers scenarios in which preference feedback or rankings over actions are used instead of traditional rewards. Due to the lack of explicit reward signals, theoretical studies of CDBs employ the concept of a \textit{von Neumann winner}—a policy that achieves at least a 50\% expected win rate when compared to any competing policy \citep{dudik2015contextual}, in lieu of seeking an optimal policy. A distinguishing characteristic of the CDB framework is that preference information is typically provided in an online manner, whereas in human preference-based learning, one often works with a pre-collected, fixed dataset containing offline preferences between action pairs \citep{yan2022human}. In a similar vein, \textit{preference-based RL} (PbRL) learns using binary preference comparisons derived from an underlying, unknown ‘scoring’ or reward function \citep{BusaFekete2014,ruiz2023dueling}. Existing PbRL approaches, including those capable of leveraging off-policy preference data, generally proceed by explicitly estimating the hidden scoring function—also referred to as the reward model—before optimizing it to improve the policy \citep{jain2013learning,BusaFekete2014,christiano2017deep,sadigh2017active,kupcsik2018learning}. By contrast, our proposed method introduces a direct policy optimization technique that learns to meet preference requirements in a single stage, without an explicit intermediate reward modelling step.

\section{Preliminaries}\label{section:prelims}

We briefly outline the RLHF procedure as described in \citeauthor{ziegler2020finetuning}, with extensions in subsequent works \citep{stiennon2022learning, bai2022training, ouyang2022training}. This approach generally follows three main steps: (1) supervised fine-tuning (SFT); (2) gathering preferences and fitting a reward model; (3) reinforcement learning based optimization.

\textbf{SFT}: The RLHF workflow commences with fine-tuning a pre-trained language model using supervised learning on curated, high-quality datasets that match the target task(s) (such as dialogue or summarization), producing the model $\pi^\text{SFT}$.

\textbf{Reward Modelling Phase}: Subsequently, prompts $x$ are fed into the SFT model to generate answer pairs $(y_1, y_2)\sim \pi^\text{SFT}(y \mid x)$. These output pairs are evaluated by human annotators, who indicate a preferred response, formalized as $y_w\succ y_l \mid x$, where $y_w$ and $y_l$ correspond to the selected and rejected completions among $(y_1, y_2)$, respectively. Although the data originate from an unknown, latent reward function $r^*(y, x)$, its exact form remains inaccessible. To characterize preferences, several modelling approaches exist—most notably, the Bradley-Terry (BT) \cite{bradley1952rankanalysis} model is frequently used (with more flexible models such as Plackett-Luce \citep{plackett1975analysis, luce2012individual} applicable if multi-way ranked outputs are available). According to the BT model, the underlying human preference probability distribution $p^*$ is expressed as:

Given a static comparison dataset $\mathcal{D} = \left\{x^{(i)}, y_w^{(i)}, y_l^{(i)}\right\}_{i=1}^N$ sampled from $p^*$, we can introduce a parameterized reward function $r_{\phi}(x, y)$ and fit its parameters through maximum likelihood estimation. By interpreting the problem as binary classification, the negative log-likelihood loss function is written as:

\begin{equation}\label{eq:reward_model}
    \mathcal{L}_R(r_{\phi}, \mathcal{D}) = -\mathbb{E}_{(x, y_w, y_l)\sim \mathcal{D}}\Big[\log \sigma\big(r_{\phi}(x, y_w) - r_{\phi}(x, y_l)\big)\Big]
\end{equation}

where the function $\sigma$ denotes the logistic sigmoid. In the language modeling setting, the architecture $r_{\phi}(x, y)$ is commonly initialized using the SFT model $\pi^\text{SFT}(y \mid x)$, and then extended with a linear projection on top of the last transformer layer to yield a scalar output representing the reward~\cite{ziegler2020finetuning}. To stabilize learning and reduce reward variance, standard practice is to center the rewards, enforcing $\mathbb{E}_{x,y\sim \mathcal{D}}[r_\phi(x, y)] = 0$ for all $x$.

\textbf{RL Fine-Tuning Phase}: At this stage, the reward function trained in the prior step serves as the feedback signal to the language model. Consistent with earlier works~\citep{jaques2017sequence, jaques2020human}, policy optimization is framed as:

\begin{equation}\label{eq:RL}
\max_{\pi_{\theta}}  \mathbb{E}_{x\sim \mathcal{D}, y\sim \pi_{\theta}(y \mid x)}\left[r_{\phi}(x, y)\right] - \beta \mathbb{D}_{\textrm{KL}}\left[\pi_{\theta}(y \mid x) \| \pi_\text{ref}(y \mid x)\right],
\end{equation}

where $\beta$ governs the strength of the regularization toward the reference policy $\pi_\text{ref}$, which is typically set as the SFT-initialized policy $\pi^\text{SFT}$. In practical implementations, the policy $\pi_\theta$ is also initialized from $\pi^\text{SFT}$. Incorporating this regularization is essential: it prevents the policy from straying excessively from the data distribution where the reward model is reliable, and helps preserve diversity in generated outputs, countering mode collapse toward only high-reward responses. Given the discrete output space of language, direct optimization of this objective is infeasible with gradient-based methods; hence, reinforcement learning algorithms are used instead. The conventional technique~\cite{ziegler2020finetuning, stiennon2022learning, bai2022training, ouyang2022training} reformulates the reward as $r(x, y) = r_{\phi}(x, y) - \beta (\log \pi_{\theta}(y \mid x) - \log \pi_\text{ref}(y \mid x))$, and maximizes it through Proximal Policy Optimization (PPO)~\cite{schulman2017proximal}.

\section{Direct Preference Optimization}\label{sec:DPO}

Addressing the scalability limitations of RL algorithms in large-scale settings like language model fine-tuning, we seek an alternative method for optimizing policies with preference data. Our framework distinguishes itself from RLHF baselines by directly exploiting preference-based supervision rather than decoupling reward model fitting and policy optimization. By employing a specific parameterization for the reward model, we can analytically recover its optimal policy, bypassing the necessity for iterative RL-based optimization. As detailed below, our central contribution is to convert the mapping from reward functions to their associated optimal policies into a reparameterization, so that the loss function can be expressed in terms of policy parameters instead of reward parameters. This change-of-variables enables us to optimize the policy network with respect to empirical preference data directly, all while preserving modeling assumptions such as the Bradley-Terry preference structure. Consequently, the policy serves dually as the language model and as an implicit reward estimator.

\textbf{Deducing the DPO objective.} We begin by adopting the reinforcement learning (RL) objective as stated in Eq.~\ref{eq:RL}, incorporating a generic reward function $r$. Drawing from established work~\citep{peters2007reinforcement, peng2019advantage, korbak2022reinforcement, go2023aligning}, one can readily verify that the optimal solution to the KL-regularized reward maximization problem described in Eq.~\ref{eq:RL} is given by:

\begin{equation}\label{eq:op_policy}
    \pi_r(y\mid x) = \frac{1}{Z(x)}\pi_\text{ref}(y\mid x)\exp\left(\frac{1}{\beta}r(x, y)\right),
\end{equation}

where the normalizing constant $Z(x) =\sum_{y}\pi_\text{ref}(y\mid x)\exp\left(\frac{1}{\beta}r(x, y)\right)$ ensures the policy sums to one. Further elaboration on this derivation can be found in Appendix \ref{app:derivation1}. Direct computation of $Z(x)$, even when the ground-truth reward $r^*$ is estimated via MLE as $r_{\phi}$, remains computationally challenging~\citep{korbak2022reinforcement, go2023aligning}. Consequently, this parameterization is often impractical for large-scale use. Nevertheless, it is possible to reformulate Eq.~\ref{eq:op_policy} to express the reward as a function of the optimal policy $\pi_r$, the reference policy $\pi_\text{ref}$, and the undetermined normalizer $Z(\cdot)$. Taking the logarithm of both sides of Eq.~\ref{eq:op_policy} and rearranging, we arrive at:

\begin{equation}\label{eq:main_eq}
    r(x,y) =\beta \log \frac{\pi_r(y\mid x)}{\pi_\text{ref}(y\mid x)} + \beta \log Z(x).
\end{equation}

This reparameterization may be applied to both the true reward $r^*$ and its associated optimal policy $\pi^*$. Importantly, in the Bradley-Terry framework, only differences in reward between completions are relevant: ${p^*(y_1 \succ y_2 \mid x) = \sigma(r^*(x, y_1) - r^*(x, y_2))}$. Substituting the above reward reparameterization (Eq.~\ref{eq:main_eq}) into the Bradley-Terry preference model (Eq.~\ref{eq:bradley-terry}), the partition function $Z(x)$ drops out. As a result, the likelihood of human preferences can be written solely in terms of $\pi^*$ and $\pi_\text{ref}$. Thus, for the Bradley-Terry model, the policy $\pi^*$ optimized via RLHF fulfills the following:

\begin{equation}\label{eq:objective}
    p^*(y_1\succ y_2 \mid x)=\frac{1}{1 + \exp\left(\beta \log \frac{\pi^*(y_2\mid x)}{\pi_\text{ref}(y_2\mid x)} - \beta \log \frac{\pi^*(y_1\mid x)}{\pi_\text{ref}(y_1\mid x)}\right)}
\end{equation}

A complete derivation is available in Appendix~\ref{app:derivation2}. While this construction is rooted in the Bradley-Terry model, analogous formulations can be derived for broader preference structures such as the Plackett-Luce models~\citep{plackett1975analysis, luce2012individual}, as discussed in Appendix~\ref{app:plackett_luce_models}.

With the human preference probabilities now parameterized by the policy rather than the reward, we may pose the problem as a maximum likelihood estimation over policy parameters $\pi_\theta$. Mirroring the methodology for reward modeling (see Eq.~\ref{eq:reward_model}), we define the DPO objective as follows:

\begin{equation}\label{eq:optimum_model}
    \mathcal{L}_\text{DPO}(\pi_{\theta}; \pi_\text{ref}) = -\mathbb{E}_{(x, y_w, y_l)\sim \mathcal{D}}\left[\log \sigma \left(\beta \log \frac{\pi_{\theta}(y_w\mid x)}{\pi_\text{ref}(y

Through this approach, an implicit reward function is learned via an alternative parameterization, for which the resulting optimal policy coincides with $\pi_\theta$. Furthermore, since our method is mathematically equivalent to training a reparametrized Bradley-Terry model, it inherits various theoretical guarantees—such as consistency—assuming appropriate conditions on the distribution of preference data, as discussed by \cite{bong2022generalized}. We elaborate further on the theoretical aspects of DPO and its connections to related work in Section~\ref{sec:theory}.

\textbf{What is the mechanism of the DPO update?} To gain a deeper mechanistic insight into DPO, it is informative to examine the gradient of the objective $\mathcal{L}_\text{DPO}$. This gradient with respect to the model parameters $\theta$ can be expressed as:

\begin{multline*}\label{eq:gradient}
    \nabla_\theta \mathcal{L}_\text{DPO}(\pi_\theta;\pi_\text{ref}) = \\ -\beta\mathbb{E}_{(x, y_w, y_l) \sim \mathcal{D}} \bigg[\underbrace{\sigma(\hat{r}_\theta(x, y_l) - \hat{r}_\theta (x, y_w))}_\text{higher weight when reward estimate is wrong}\bigg[\underbrace{\nabla_\theta\log \pi(y_w \mid x)}_\text{increase likelihood of $y_w$} - \underbrace{\nabla_\theta\log\pi(y_l \mid x)}_\text{decrease likelihood of $y_l$}\bigg]\bigg],
\end{multline*}

where the implicitly defined reward is given by $\hat{r}_\theta(x, y) = \beta \log \frac{\pi_\theta(y \mid x)}{\pi_\text{ref}(y \mid x)}$, utilizing both the parameterized language model $\pi_\theta$ and the reference policy $\pi_\text{ref}$ (see Section~\ref{sec:theory} for additional detail). Intuitively, this gradient functions to increase the model’s probability for more favored completions $y_w$, while simultaneously suppressing the probabilities assigned to less preferred completions $y_l$. Notably, each training sample is assigned a weight reflecting the extent to which the implicit reward model $\hat{r}_\theta$ incorrectly scores the dispreferred completions above the preferred ones, modulated by the parameter $\beta$. This captures both the severity of misordering by the implicit reward model and the impact of the KL regularization. Empirically, we find that this dynamic weighting is critical: ablations without the weighting term, which naïvely omit this coefficient, tend to lead the language model to collapse or degenerate (refer to Appendix Table~\ref{tab:unlikelihood_generations}).

\textbf{DPO summary.} The typical workflow of DPO consists of the following steps: 1) For each prompt $x$, generate sample completions $y_1, y_2 \sim \pi_\text{ref}(\cdot \mid x)$, then annotate these samples with human preference signals to assemble an offline preference dataset $\mathcal{D} = \{x^{(i)}, y_w^{(i)}, y_l^{(i)}\}_{i=1}^N$; 2) Train the language model $\pi_\theta$ to minimize the DPO loss $\mathcal{L}_\text{DPO}$ using the specified $\pi_\text{ref}$, dataset $\mathcal{D}$, and a chosen $\beta$. 
Practically, it is preferable to leverage publicly available preference datasets instead of producing new completions and collecting fresh human preference labels. When these datasets have been generated using a supervised fine-tuned model $\pi^\text{SFT}$, we set $\pi_\text{ref} = \pi^\text{SFT}$ if possible. In the absence of $\pi^\text{SFT}$, we fit $\pi_\text{ref}$ by maximizing the likelihood of the preferred completions ${(x, y_w)}$, formally: ${\pi_\text{ref} = \argmax_{\pi}\mathbb{E}_{x, y_w \sim \mathcal{D}}\left[\log \pi(y_w \mid x)\right]}$. This step helps address potential distribution mismatches between the (inaccessible) true reference distribution and the practical reference model $\pi_\text{ref}$ deployed within DPO. Additional implementation details and hyperparameter choices can be found in Appendix~\ref{app:implementation}.

\section{Theoretical Analysis of DPO}
This section develops theoretical insights into the DPO framework, clarifies its properties, and situates its benefits relative to actor-critic methods employed in RLHF (e.g., PPO~\cite{schulman2017proximal}), identifying where DPO offers improvements.

\label{sec:theory}

\subsection{Your Language Model Is Secretly a Reward Model}
DPO manages to avoid the explicit learning of a reward function and the subsequent policy optimization through reinforcement learning by employing a maximum likelihood-based loss. Notably, its optimization formulation, as presented in Eq. \ref{eq:main_eq}, corresponds to a Bradley-Terry model with a reward defined as $r^*(x, y) = \beta \log\frac{\pi^*_\theta(y \mid x)}{\pi_\text{ref}(y \mid x)}$. The parameterized policy model $\pi_{\theta}$ is then optimized analogously to reward model training described in Eq. \ref{eq:reward_model}, once the reparameterization is applied. In what follows, we systematically establish the theoretical foundation for this alternative parametrization, demonstrate that it is as expressive as traditional reward modeling, and verify that it permits precise retrieval of the optimal policy. To start, we introduce an equivalence concept for reward functions.

\begin{definition}
We define two reward functions $r(x, y)$ and $r'(x, y)$ as equivalent if and only if ${r(x, y) - r'(x, y) = f(x)}$ holds for some function $f$.     
\end{definition}

This relation clearly satisfies the properties of an equivalence relation, thereby dividing the set of all reward functions into distinct equivalence classes. This allows us to formalize two key lemmas:

\begin{lemma}\label{lemma:same_prefrence} Within the Plackett-Luce family—encompassing the Bradley-Terry model—any pair of reward functions belonging to the same equivalence class will yield identical preference distributions.
\end{lemma}

\begin{lemma}\label{lemma:same_policy}
    For the constrained RL problem, optimal policies derived from any two reward functions in the same equivalence class will be identical.
\end{lemma}

Since the proofs are immediate, we postpone them to Appendix \ref{app:lemma1}. The first lemma articulates a well-documented identifiability issue inherent in the Plackett-Luce framework \cite{plackett1975analysis}. Because of this ambiguity, it becomes necessary to enforce additional identifiability conditions in order to derive meaningful MLE solutions from Eq. \ref{eq:reward_model} \cite{bong2022generalized}. The second lemma indicates that every member within an equivalence class leads to the same optimal policy, and, consequently, recovering any representative from the class suffices for our objectives. In Appendix~\ref{app:thm1}, we establish the following theorem:

\begin{theorem}\label{thm:main}
    Provided mild regularity conditions, every reward equivalence class consistent with the Plackett-Luce models (including Bradley-Terry) can be captured via the reparameterization ${r(x, y) = \beta \log \frac{\pi(y\mid x)}{\pi_\text{ref}(y\mid x)}}$ for some choice of model $\pi(y\mid x)$ and fixed reference distribution $\pi_\text{ref}(y \mid x)$.
\end{theorem}

\begin{sproof}
    Take an arbitrary reward function $r(x, y)$, which defines its own optimal policy $\pi_r(y \mid x)$ as described in Eq. \ref{eq:op_policy}. We aim to demonstrate that an equivalent reward function to $r$ can be constructed using the proposed reparameterization. To this end, introduce the mapping $f$:
\begin{equation}
    f(r; \pi_\text{ref}, \beta)(x, y) = r(x, y) - \beta\log\sum_{y}\pi_\text{ref}(y\mid x)\exp\left(\frac{1}{\beta}r(x, y)\right)
\end{equation}
Here, $f$ acts by normalizing $r(x, y)$ through the logarithm of the corresponding partition function for $\pi_r$. The normalization term only depends on the input $x$, thus ensuring $f(r; \pi_\text{ref}, \beta)(x, y)$ remains within the equivalence class of the original $r(x, y)$. Substituting the reward structure from Eq.~\ref{eq:main_eq}, which applies to any $r$, yields $f(r; \pi_\text{ref}, \beta)(x, y) = \beta \log \frac{\pi_r(y\mid x)}{\pi_\text{ref}(y\mid x)}$. This shows that $f$ produces an equivalence-class representative in the specified reparameterized form, confirming that the proposed model loses no generality.
\end{sproof}

An alternative perspective on Theorem~\ref{thm:main} is that it precisely identifies which reward function, within any given equivalence class, is selected by the DPO reparameterization: specifically, the reward function that fulfills

\begin{equation}\label{eq:lag_p}
     \sum_{y}\underbrace{\pi_\text{ref}(y\mid x)\exp\left(\frac{1}{\beta}r(x, y)\right)}_{=\pi(y\mid x)\text{, using Thm.~\ref{thm:main} reparam.}} = 1,
\end{equation}

In other words, this ensures $\pi(y\mid x)$ constitutes a legitimate probability distribution (i.e., nonnegative probabilities summing to one). Notably, drawing on Eq.~\ref{eq:op_policy}, Eq.~\ref{eq:lag_p} corresponds to the partition function associated with the optimal policy induced by the reward $r(x, y)$. The fundamental idea underlying the DPO algorithm is that, by enforcing particular constraints on the underdetermined Plackett-Luce (and, as a specific case, Bradley-Terry) family of preference models, it is possible to retain the expressive capacity of the set of reward models, while at the same time ensuring that the optimal policy from Eq.~\ref{eq:op_policy} remains analytically solvable for any prompt $x$.

\subsection{Instability of Actor-Critic Algorithms} The presented framework can also be leveraged to identify sources of instability in common actor-critic methods used for RLHF, such as PPO. Following the RLHF methodology, we focus on the RL fine-tuning phase as detailed in Section \ref{section:prelims}. Connections can be made to the control-as-inference paradigm \cite{levine2018reinforcement} for the constrained RL scenario depicted by \ref{eq:RL}. Assuming a parameterized policy model $\pi_{\theta}(y\mid x)$, the objective is to minimize $\mathbb{D}_{\text{KL}}[\pi_{\theta}(y|x) \mid \mid \pi^*(y\mid x)]$, where $\pi^*$ is the optimal policy derived from Eq. \ref{eq:optimum_model}, itself induced by the reward function $r_{\phi}(y, x)$. Through straightforward manipulations, this leads to the following objective:

\begin{equation}\label{eq:AC}
    \max_{\pi_{\theta}}\mathbb{E}_{\pi_{\theta}(y\mid x)}\bigg[\underbrace{r_{\phi}(x, y) -\beta\log\sum_{y}\pi_\text{ref}(y\mid x)\exp\left(\frac{1}{\beta}r_{\phi}(x, y)\right)}_{f(r_{\phi}, \pi_\text{ref}, \beta)} - \underbrace{\beta\log\frac{\pi_{\theta}(y\mid x)}{\pi_\text{ref}(y\mid x)}}_{\text{KL}}\bigg]
\end{equation}

The objective being optimized here is identical to that employed in earlier studies
\citep{ziegler2020finetuning, stiennon2022learning, bai2022training, ouyang2022training}, which use the DPO-equivalent reward formulation within the $r_{\phi}$ reward family. Within this framework, the normalization component in $f(r_{\phi}, \pi_\text{ref}, \beta)$ can be understood as the soft value function associated with the reference policy $\pi_\text{ref}$. Though this normalization does not alter the optimal policy, its absence can lead to high-variance policy gradient estimates, potentially destabilizing the learning process. To manage this term, one option is to introduce a learned value function, although this approach can be challenging to optimize in practice. Another common strategy, found in preceding work, is to normalize rewards based on a human completion baseline, effectively yielding a one-sample Monte Carlo approximation to the normalization factor. The DPO reparameterization, in contrast, produces a reward function that entirely circumvents the need for such baselines.

\section{Experiments}
Here, we conduct empirical assessments of DPO’s capacity to train policies exclusively using preference data. We begin by examining a controlled text-generation scenario to address the question: relative to prevalent preference optimization methods like PPO, how effectively does DPO balance reward maximization against KL-divergence regularization with the reference policy? Subsequently, we test DPO on more complex, large-scale problems—including summarization and dialogue—within RLHF contexts. Our experiments show that, with minimal hyperparameter tuning, DPO generally matches or surpasses established baselines such as PPO-based RLHF, as well as outperforming the selection of the top trajectory from $N$ samples based on a learned reward model. Details on the experimental protocol are described below, with further information available in Appendix~\ref{app:exp_details}.

\textbf{Tasks.} Our study investigates three open-ended text generation tasks. For all tasks, methods are trained to optimize a policy using a preference dataset $\mathcal{D}=\bigl\{x^{(i)}, y_w^{(i)}, y_l^{(i)}\bigr\}_{i=1}^N$. In the \textbf{controlled sentiment generation} task, $x$ serves as the initial segment of an IMDb movie review \cite{maas-EtAl:2011:ACL-HLT2011}, with the objective for the policy to produce a continuation $y$ demonstrating positive sentiment. To standardize evaluation, this experiment employs automatically generated preference pairs via a pre-trained sentiment classifier, selecting pairs such that $p(\text{positive}\mid x,y_w)>p(\text{positive}\mid x,y_l)$. The SFT approach entails fine-tuning GPT-2-large to convergence on the training partition of the IMDb dataset (additional specifics in App~\ref{app:sentiment_details}). In the \textbf{summarization} task, $x$ corresponds to a Reddit forum post, and the policy is expected to generate a summary $y$ encapsulating the main ideas. Consistent with established practice, we utilize the Reddit TL;DR dataset \citep{volske-etal-2017-tl} supplemented by human preference data collected by \citeauthor{stiennon2022learning}. For SFT, we deploy a model fine-tuned on human-composed post summaries\footnote{\url{https://huggingface.co/CarperAI/openai_summarize_tldr_sft}} using the TRLX \citep{leandro_von_werra_2023_7790115} library for RLHF. These human preferences were annotated by \citeauthor{stiennon2022learning} from outputs of a distinct, comparably-trained SFT model. The \textbf{single-turn dialogue} task features $x$ as a human-issued prompt, which can span diverse topics including astrophysics questions or requests for relationship guidance. Here, the policy is tasked with generating an informative and engaging reply $y$. This setting leverages the Anthropic Helpful and Harmless dialogue corpus \citep{bai2022training}, which comprises 170k human–assistant interactions. Each sample concludes with a pair of model-generated replies and a human judgment indicating the preferred answer. Due to the lack of an existing SFT checkpoint in this domain, we construct the SFT model by fine-tuning a general-purpose language model on the set of preferred responses.

\textbf{Evaluation.} Our empirical studies employ two distinct strategies for evaluation. To systematically examine how well each algorithm optimizes the constrained reward maximization objective, we measure their performance in the controlled sentiment generation scenario by mapping the frontier formed by reward achieved versus KL-divergence from the reference policy. This assessment is feasible because access to the true reward function (provided by a sentiment classifier) is available. Conversely, in realistic scenarios where the exact reward function is unknown, we rely on the \textit{win rate} metric against a baseline policy. For this purpose, GPT-4 serves as an automated evaluator, approximating human judgment of summary quality and response utility within summarization and single-turn dialogue tasks, respectively. In the summarization case, the test set's reference summaries function as the baseline; in dialogue, the chosen responses marked as preferred within the dataset act as the baseline. While prior work indicates that language models may surpass traditional metrics as evaluators \citep{Chen2023ExploringTU}, we further support the validity of GPT-4 evaluation by conducting a human study as detailed in Sec.~\ref{sec:human-judgments}. The study reveals that GPT-4's assessments exhibit high concordance with human judgments, often matching or exceeding the level of agreement observed between different human annotators.

\textbf{Methods.} Besides DPO, our evaluation encompasses a variety of established methods for aligning language models with human preferences. At the simplest level, we include zero-shot prompting with \textbf{GPT-J} \citep{gpt-j} for summarization and 2-shot prompting with \textbf{Pythia-2.8B} \citep{biderman2023pythia} for the dialogue task. The \textbf{SFT} model and \textbf{Preferred-FT}—the latter fine-tuned via supervised learning on the selected completion $y_w$ (generated by the SFT model in sentiment and summarization settings or by a generic language model in dialogue)—are also evaluated. We incorporate the \textbf{Unlikelihood}~\citep{welleck2019neural} technique as another pseudo-supervised baseline, where the policy is explicitly optimized to increase the likelihood of $y_w$ and decrease the likelihood of $y_l$, leveraging a tunable coefficient $\alpha\in[0,1]$ for the ‘unlikelihood’ term. Moreover, we assess \textbf{PPO} \citep{schulman2017proximal}, which operates with a reward model trained on preference data, and \textbf{PPO-GT}, an oracle variant utilizing the ground-truth reward in the controlled sentiment setting. In the sentiment experiments, PPO-GT is instantiated both as a standard implementation \cite{leandro_von_werra_2023_7790115} and in an enhanced form with reward normalization and fine-tuned hyperparameters, the latter improvements also being applied to PPO with learned rewards. As an additional baseline, we consider \textbf{Best of $N$}, which involves sampling $N$ responses from either the SFT model (or Preferred-FT in dialogue) and selecting the top response as ranked by a reward model trained on preferences. Although highly performant, this method decouples the reward model from PPO optimization but suffers from inefficiency, as it requires generating $N$ candidates for every input during evaluation, rendering it impractical for larger $N$.

\subsection{How well can DPO optimize the RLHF objective?}

In standard RLHF approaches, the objective of maximizing the KL-constrained reward requires striking a balance between exploiting the reward function and ensuring the policy does not stray significantly from the reference. Consequently, comparisons across methods must jointly consider the achieved reward and the corresponding KL divergence; obtaining a marginally higher reward with a substantial increase in KL may be suboptimal. Figure~\ref{fig:frontier-tldr-main} illustrates the tradeoff between reward and KL divergence—the frontier—for different methods within the sentiment analysis task. For each method, we conduct several independent training sessions, each utilizing distinct settings for policy conservativeness (target KL $\in\{3,6,9,12\}$ for PPO, $\beta \in \{0.05,0.1,1,5\}$ for DPO, $\alpha\in\{0.05,0.1,0.5,1\}$ for unlikelihood, and various random seeds for preferred-FT), resulting in a total of 22 experimental runs. Every 100 training steps, up to convergence, we evaluate each resulting policy on a standardized set of test prompts, calculating both the mean reward with respect to the ground truth reward and the average sequence-level KL\footnote{Specifically, the aggregate of KL-divergence values computed at each timestep.} relative to the reference policy $\text{KL}\left(\pi\mid \mid \pi_\text{ref}\right)$. Our analysis reveals that DPO establishes a markedly superior efficiency frontier: it attains higher rewards without incurring excessive KL divergence. This finding stands out for several reasons. Notably, although both DPO and PPO aim to optimize the same objective, DPO does so with substantially greater efficiency, consistently outperforming PPO in the reward/KL tradeoff. Moreover, DPO surpasses even PPO when PPO has access to the true rewards (PPO-GT), underscoring DPO’s robust performance.

\subsection{Can DPO scale to real preference datasets?}
\label{sec:dpo-real-datasets}

We proceed to assess DPO's fine-tuning capabilities on real-world preference tasks involving summarization and single-turn dialogue. In the domain of summarization, established automated metrics like ROUGE have been demonstrated to correlate weakly with human judgments~\citep{stiennon2022learning}. Previous studies have shown that leveraging PPO for language model fine-tuning, using human preference data, can yield superior summaries. Our comparative analysis involves generating model completions on the test partition of the TL;DR summarization dataset and measuring the mean win rate relative to reference completions from the same test set. Each method's outputs are sampled across temperature settings from 0.0 up to 1.0; the resulting win rates are displayed in Figure~\ref{fig:frontier-tldr-main} (right). All methods—including DPO, PPO, and Preferred-FT—are fine-tuned on a common GPT-J SFT model\footnote{\url{https://huggingface.co/CarperAI/openai_summarize_tldr_sft}}. Our results indicate that DPO attains a win rate close to 61\% at a sampling temperature of 0.0, surpassing PPO, which reaches approximately 57\% at its best performance, also at temperature 0.0. Furthermore, DPO achieves a higher peak win rate than the best of $N$ baseline approach. It is important to highlight that minimal tuning was performed for DPO’s $\beta$ hyperparameter; thus, there is likely additional untapped performance. Notably, DPO exhibits greater stability across temperature settings, whereas PPO's effectiveness decreases, approaching that of the base GPT-J model at elevated temperatures. The Preferred-FT method, by contrast, shows only marginal improvement over the base SFT model. Finally, Section~\ref{sec:human-judgments} presents a direct comparison between DPO and PPO via human assessment: DPO completions at temperature 0.25 were selected as preferable in 58\% of trials over PPO outputs at the same temperature.

For single-turn dialogue evaluation, we focus on the portion of the Anthropic HH dataset \citep{bai2022training} test set containing a single human-assistant exchange. To determine win rates for each method, GPT-4 assessments reference the preferred completions within the test set. In the absence of a standard SFT model for this domain, we initialize with a pre-trained Pythia-2.8B, apply Preferred-FT to develop a reference model tuned to produce in-distribution completions, and subsequently fine-tune it using DPO. Additionally, we include as baselines both the best of 128 Preferred-FT completions (noting that the Best of $N$ benchmark saturates at $N = 128$ for this task, as illustrated in Appendix Figure~\ref{fig:best-of-n}) and a 2-shot prompt configuration for the Pythia-2.8B base model. Our findings show that DPO matches or exceeds the best outcomes across the optimal sampling temperatures for each approach. Moreover, we assess a PPO-based RLHF model, trained on the Anthropic HH dataset \footnote{\url{https://huggingface.co/reciprocate/ppo_hh_pythia-6B}} from a reputable implementation \footnote{\url{https://github.com/CarperAI/trlx/tree/main/examples/hh}}, but we are unable to identify a prompting strategy or temperature that surpasses the base Pythia-2.8B’s performance. Considering our TL;DR task results and the fact that both Best of 128 and PPO methods target the same reward signal, we use Best of 128 as an approximate indicator for PPO-level results. In summary, DPO stands out as the sole computationally efficient approach that reliably surpasses the preferred completions in the Anthropic HH test set, offering comparable or superior results to the resource-intensive Best of 128 standard. Lastly, Figure~\ref{fig:dialogue-main} demonstrates that DPO attains peak performance after relatively few training steps.

\subsection{Generalization to a new input distribution}

To further assess the robustness of PPO and DPO models under distributional changes, we apply the policies learned from our Reddit TL;DR summarization study to a distinct domain, specifically to news articles within the test set of the CNN/DailyMail dataset \citep{nallapati-etal-2016-abstractive}. The evaluation uses the optimal sampling temperatures identified in the TL;DR task (0 and 0.25). The findings, detailed in Table~\ref{tab:ood}, include GPT-4 win rates based on comparisons to the ground-truth news summaries, utilizing the same GPT-4 (C) prompt as for Reddit TL;DR, but with the term ``forum post'' modified to ``news article''. On this unfamiliar input distribution, DPO consistently exceeds PPO in performance by a notable amount. These results offer preliminary support that DPO-trained policies demonstrate generalization capabilities comparable to those of PPO, despite DPO not leveraging the extra unlabeled Reddit TL;DR prompt data accessible to PPO.

\subsection{Validating GPT-4 judgments with human judgments}
\label{sec:human-judgments}
A human evaluation is carried out to assess the trustworthiness of GPT-4’s assessments, utilizing the TL;DR summarization experiment outcomes along with two variants of GPT-4 prompts. The \textbf{GPT-4 (S)} (simple) prompt asks for a judgment about which summary most effectively captures the essential information in the post. The \textbf{GPT-4 (C)} (concise) prompt requests a choice based on informativeness as well as conciseness, and is included because, under the \textbf{GPT-4 (S)} prompt, GPT-4 often shows a preference for longer and more redundant summaries, unlike human evaluators. Full prompt texts are in Appendix~\ref{app:prompts}. Three system comparisons are evaluated: the top-performing model (DPO, temp. 0.25), the lowest-performing (PPO, temp. 1.0), and an intermediate baseline (SFT, temp. 0.25). Each of these is benchmarked against greedily decoded PPO outputs (using its optimal temperature), aiming to span a range of sample qualities. The results show that GPT-4, under both prompt conditions, aligns with human judgment as frequently as humans agree with each other. This supports the use of GPT-4 as a valid stand-in for human evaluators (note that multiple human ratings were gathered only for the DPO and PPO-1 matchups, given the limited pool of human participants). In general, the \textbf{GPT-4 (C)} prompt yields win rates more closely mirroring human preferences, and so it is adopted for the principal analyses in Section~\ref{sec:dpo-real-datasets}. Further details on the human evaluation setup, including screenshots of the web-based rater interface and the roster of human raters, are given in Appendix~\ref{app:human-study}.

\section{Discussion}
Preference-based learning offers a robust and scalable approach for developing language models that are both capable and well-aligned with human values. In this work, we presented DPO, a streamlined framework that enables language models to learn from preferences without relying on reinforcement learning methodologies. Instead of reshaping the preference learning challenge to fit conventional RL frameworks and their accompanying algorithms, DPO establishes a connection between language model policies and reward structures, which makes it possible to optimize models to meet human preferences \textit{directly} using a basic cross-entropy loss function—eliminating the need for reinforcement learning and without compromising generality. Remarkably, DPO achieves comparable or superior results to current RLHF approaches, such as those utilizing PPO, with little requirement for hyperparameter adjustment. This substantially simplifies and broadens the accessibility of training language models based on human feedback.

\textbf{Limitations \& Future Work.} The findings of this study surface several avenues for further investigation. One open question concerns the out-of-distribution generalization of policies learned with DPO in comparison to those trained with an explicit reward signal. While our preliminary results indicate that DPO-trained policies generalize on par with models based on PPO, more systematic evaluation is warranted. For instance, it remains to be seen whether methods such as self-labeling using the DPO policy can leverage unlabeled prompts as effectively. Additionally, the nature of reward over-optimization within the direct preference optimization paradigm merits further analysis, particularly in light of the modest performance drop observed in Figure~\ref{fig:dialogue-main}-right—potentially a manifestation of such phenomena. Although our current evaluation encompasses models up to 6B parameters, scaling DPO to cutting-edge architectures that are orders of magnitude larger remains an intriguing and important area for subsequent research. In the context of model assessment, our experiments show that GPT-4's computed win rates are sensitive to prompt wording; future studies may focus on refining strategies to obtain reliable quality judgments from automated evaluation systems. Beyond language modeling, DPO presents promising opportunities for application in training generative models across different domains.